\documentclass[11pt,a4paper]{article}

% ===== Codificación e idioma =====
\usepackage[utf8]{inputenc}
\usepackage[T1]{fontenc}
\usepackage[spanish,es-tabla]{babel}
\usepackage{lmodern}
\usepackage{microtype}

% ===== Geometría de página =====
\usepackage[a4paper, top=2.5cm, bottom=2.5cm, left=3cm, right=2.5cm]{geometry}
\usepackage{setspace}
\onehalfspacing
\usepackage{amsmath}

% ===== Gráficos y tablas =====
\usepackage{graphicx}
\usepackage{float}
\usepackage{tabularx}
\usepackage{booktabs}
\usepackage{array}
\usepackage{multirow}
\usepackage{caption}
\captionsetup{font=small,labelfont=bf,format=plain}
\usepackage{subcaption}
\usepackage[table]{xcolor}
\usepackage{enumitem}

% ===== Enlaces y referencias =====
\usepackage{hyperref}
\hypersetup{
  colorlinks=true,
  linkcolor=blue!80!black,
  citecolor=blue!80!black,
  urlcolor=blue!80!black
}

% ===== Encabezado y pie =====
\usepackage{fancyhdr}
\pagestyle{fancy}
\fancyhf{}
\fancyhead[L]{\small UNSAAC -- IF614 -- Ing. de Software I}
\fancyhead[R]{\small CGB Academy -- Informe de Software}
\fancyfoot[C]{\thepage}
\renewcommand{\headrulewidth}{0.4pt}
\renewcommand{\footrulewidth}{0.4pt}

\begin{document}

% ===== PORTADA =====
\begin{titlepage}
  \centering
  {\scshape\LARGE Universidad Nacional de San Antonio Abad del Cusco \par}
  \vspace{0.3cm}
  {\scshape\large Facultad de Ingeniería Eléctrica, Electrónica, Informática y Mecánica\par}
  \vspace{0.2cm}
  {\scshape\normalsize Escuela Profesional de Ingeniería de Sistemas e Informática\par}
  \vspace{1.2cm}

  {\Large\textbf{CURSO: INGENIERÍA DE SOFTWARE I (IF614)}\par}
  \vspace{0.8cm}
  
  \rule{\linewidth}{0.5mm}\\[0.4cm]
  {\huge\bfseries Centro de Capacitación e Inducción\\[0.2cm] CGB Academy\par}
  \vspace{0.3cm}
  {\large\textit{Informe de Ingeniería de Software: Requisitos, Arquitectura y Método Híbrido}\par}
  \rule{\linewidth}{0.5mm}\\[1.2cm]

  \textbf{Entidad Beneficiaria:}\\[0.2cm]
  CGB Academy (CIIP LATAM, GEOMINA LATAM y BIOMEDIC)\\[0.8cm]

  \textbf{Integrantes del Equipo de Desarrollo:}\\[0.3cm]
  \begin{tabular}{ll}
    \textbf{Barazorda Cuellar, Hector} & \textit{Scrum Master / Dev Full-Stack} \\
    \textbf{Carpio Hermoza, Alex} & \textit{Desarrollador Full-Stack} \\
    \textbf{Quispe Mamani, Domingo de Guzman} & \textit{Desarrollador Full-Stack} \\
    \textbf{Saire Bustamante, Edmil Jampier} & \textit{Product Owner / Dev Full-Stack} \\
    \textbf{Ticona Jancco, Ronaldo} & \textit{Desarrollador Full-Stack} \\
  \end{tabular}
  \vfill
  {\large Cusco, Perú\par}
  {\large 2026\par}
\end{titlepage}

\tableofcontents
\newpage

\section{Datos Generales del Proyecto}
\begin{itemize}[leftmargin=*]
    \item \textbf{Denominación del Sistema:} Centro de Capacitación e Inducción CGB Academy.
    \item \textbf{Organización Beneficiaria:} CGB Academy, integrando los programas de CIIP LATAM, GEOMINA LATAM y BIOMEDIC sin generar plataformas dispersas.
    \item \textbf{Repositorio Central en GitHub:} \url{https://github.com/milith0kun/proyecto-ing-software}
    \item \textbf{Ramas de Trabajo:} \texttt{main} (código estable de producción) y \texttt{develop} (integración continua colaborativa).
    \item \textbf{Alcance del MVP:} 15 Historias de Usuario funcionales agrupadas en 3 Sprints de dos semanas de duración cada uno.
\end{itemize}

\section{Objetivo General y Específicos}
\subsection{Objetivo General}
Diseñar, implementar y validar el \textbf{Centro de Capacitación e Inducción CGB Academy}, una solución web modular orientada a la capacitación de colaboradores internos y a la orientación pública de estudiantes y docentes, estructurada en diapositivas interactivas (\textit{slides}), rutas de aprendizaje organizacional jerárquicas y microtests formativos, prescindiendo de la complejidad operativa de un LMS tradicional.

\subsection{Objetivos Específicos}
\begin{enumerate}[leftmargin=*]
    \item Especificar y priorizar 15 historias de usuario funcionales para el MVP con criterios de aceptación verificables en formato BDD (\textit{Dado que / Cuando / Entonces}).
    \item Diseñar una arquitectura web ligera, desacoplada y responsiva apta para desplegarse sobre una infraestructura VPS objetivo de 2 vCPU y 8 GB de RAM.
    \item Implementar el motor de creación y publicación de capacitaciones mediante \textit{slides} administrables sin requerir conocimientos de programación (\textit{no-code}).
    \item Modelar el módulo organizacional de onboarding relacionando \textit{Área $\rightarrow$ Puesto $\rightarrow$ Ruta $\rightarrow$ Capacitación}, garantizando activación segura y registro automático de progreso.
    \item Incorporar microtests formativos con retroalimentación inmediata orientados al aprendizaje sin penalizaciones punitivas.
    \item Aplicar un método ágil híbrido combinando la cadencia de Scrum, la visibilidad de Kanban y las disciplinas de ingeniería de XP.
\end{enumerate}

\section{Alcance del Sistema}
\subsection{Funcionalidades Incluidas (In-Scope -- MVP)}
\begin{itemize}[leftmargin=*]
    \item \textbf{Sprint 1 (Motor y Acceso Público):} Autenticación con control de acceso por roles (Admin/Colaborador), creador de capacitaciones, editor interactivo de slides (texto, imágenes, listas, indicaciones, botones), previsualización, gestión de estados (Borrador/Publicada) y catálogo público libre sin autenticación.
    \item \textbf{Sprint 2 (Onboarding Interno):} Mantenimiento de áreas y puestos con integridad referencial, construcción de rutas de onboarding, registro de colaboradores, activación mediante token seguro de uso único y módulo privado ``Mi Onboarding''.
    \item \textbf{Sprint 3 (Microtests y Supervisión):} Creación de evaluaciones cortas de opción múltiple, resolución con retroalimentación inmediata al colaborador y panel administrativo para supervisar el avance individual y colectivo.
\end{itemize}

\subsection{Exclusiones Expresas (Out-of-Scope)}
Se excluyen deliberadamente para el MVP: venta de cursos, carritos de compra y pasarelas de pago; aulas virtuales en vivo; emisión automatizada de diplomas o certificados oficiales; almacenamiento de video pesado o biblioteca de descargas masivas; soporte multi-idioma (versión 1.0 exclusivamente en español); y migración de historiales legados.

\section{Justificación del Proyecto}
\subsection{Justificación Práctica y Organizacional}
CGB Academy agrupa la formación especializada de CIIP LATAM, GEOMINA LATAM y BIOMEDIC. Previamente, la inducción de personal dependía de documentos manuales en PDF y explicaciones repetitivas no trazables. La plataforma automatiza la inducción, estandariza procesos corporativos y ofrece a la comunidad externa un centro de consulta claro y accesible en todo momento.

\subsection{Justificación Académica y de Ingeniería}
El proyecto aplica de manera rigurosa los estándares de la ingeniería de software moderna \cite{sommerville2011, pressman2010}: elicitación formal de requisitos mediante historias de usuario atómicas, diseño de arquitecturas eficientes para recursos limitados, control de versiones bajo GitFlow y aseguramiento continuo de la calidad.

\section{Contexto del Proyecto y Restricciones}
\subsection{Contexto Operativo}
El sistema prestará servicio a nivel nacional e internacional vía web. Los colaboradores realizarán su onboarding en horarios laborales desde sus computadoras de trabajo o dispositivos móviles, mientras que el público externo consultará libremente contenidos informativos sin fricción de registro.

\subsection{Restricciones del Sistema (Requerimientos No Funcionales)}
\begin{itemize}[leftmargin=*]
    \item \textbf{RNF-001 (Diseño Responsive):} Adaptación visual a computadoras, tablets y móviles.
    \item \textbf{RNF-002 (Idioma):} Interfaz disponible únicamente en español para la v1.0.
    \item \textbf{RNF-003 (Infraestructura):} Capacidad de operar fluidamente en VPS de 2 vCPU y 8 GB de RAM.
    \item \textbf{RNF-004 (Contenido Ligero):} Priorización de texto e imágenes optimizadas; sin videos pesados.
    \item \textbf{RNF-005 (Seguridad Interna):} Segmentación estricta de rutas privadas con redirección a login.
    \item \textbf{RNF-006 (Protección de Credenciales):} Contraseñas cifradas mediante algoritmos robustos y tokens de activación de un solo uso.
    \item \textbf{RNF-007 (Minimización de Datos):} Registro sin solicitud de datos personales sensibles no indispensables.
    \item \textbf{RNF-008 (Consistencia Visual):} Sistema de componentes estandarizado con la identidad de CGB Academy.
    \item \textbf{RNF-009 (Facilidad Administrativa):} Administración de slides mediante interfaz gráfica intuitiva (\textit{no-code}).
    \item \textbf{RNF-010 (Integridad de Información):} Protección contra eliminación huérfana de áreas o puestos con personal activo.
\end{itemize}

\section{Tecnologías Utilizadas}
\begin{table}[H]
\centering
\small
\begin{tabularx}{\linewidth}{|l|l|X|}
\hline
\rowcolor{gray!15}
\textbf{Capa / Componente} & \textbf{Tecnología} & \textbf{Justificación Técnica} \\
\hline
Base de Datos & MongoDB Atlas & Base de datos NoSQL documental en la nube, flexible y óptima para estructuras dinámicas de slides y progreso de onboarding. \\
\hline
ORM / Acceso a Datos & Prisma ORM & Modelado declarativo con conector MongoDB (\texttt{provider = "mongodb"}), tipado estático seguro y migraciones eficientes. \\
\hline
Arquitectura Full-Stack & Next.js (App Router) & Framework unificado que resuelve el Frontend (React con SSR/CSR) y el Backend (Route Handlers en \texttt{/app/api/...} y Server Actions). \\
\hline
Aplicación Web & Next.js Web App & Aplicación responsiva, óptima para SEO y tiempos de carga instantáneos en accesos públicos y privados. \\
\hline
Estilos UI & TailwindCSS / CSS Moderno & Sistema de diseño estandarizado basado en utilidades, liviano y consistente con la identidad de CGB Academy. \\
\hline
Lenguaje & TypeScript & Tipado estático integral de extremo a extremo (base de datos, endpoints y componentes de UI). \\
\hline
Control de Versiones & Git \& GitHub & Ramas: principal (\texttt{edmil-saire}), desarrollo (\texttt{desarrollo}) e individuales por desarrollador. \\
\hline
Gestión de Proyecto & Jira / GitHub Projects & Tableros Kanban por fases, sprints, estimación con Planning Poker y métricas de flujo. \\
\hline
\end{tabularx}
\caption{Pila tecnológica seleccionada para CGB Academy.}
\label{tab:tecnologias}
\end{table}

\section{Stakeholders (Matriz de Interesados)}
\begin{table}[H]
\centering
\small
\begin{tabularx}{\linewidth}{|l|l|c|X|}
\hline
\rowcolor{gray!15}
\textbf{Actor / Interesado} & \textbf{Tipo} & \textbf{Influencia} & \textbf{Expectativa Principal} \\
\hline
Colaborador Interno & Directo & Alta & Experiencia de onboarding clara, interactiva y sin fricciones técnicas. \\
\hline
Administrador & Directo & Alta & Creación sencilla de slides, asignación de rutas y visualización de progreso. \\
\hline
Público General & Externo & Media & Consulta rápida y abierta de contenidos de inducción sin requerir registro. \\
\hline
Directorio CGB & Clave & Crítica & Reducción de costos de inducción, cumplimiento del cronograma y estabilidad. \\
\hline
Equipo de Desarrollo & Interno & Alta & Código limpio, arquitectura mantenible y entregas incrementales verificadas. \\
\hline
Docente Evaluador & Académico & Crítica & Rigor metodológico en ingeniería de software y software completamente operativo. \\
\hline
\end{tabularx}
\caption{Matriz de stakeholders del proyecto.}
\label{tab:stakeholders}
\end{table}

\section{Enfoque de Desarrollo Seleccionado y Justificación}
Se ha adoptado un \textbf{Método Ágil Híbrido} \cite{sommerville2011, pressman2010}:
\begin{itemize}[leftmargin=*]
    \item \textbf{Descarte de Cascada:} Se descartó por su inviabilidad ante ajustes dinámicos de interfaz y requerimientos en sprints de corta duración.
    \item \textbf{Estructura Base Scrum \cite{scrumguide2020}:} Provee la cadencia de trabajo en sprints de 2 semanas, ceremonias de planificación, Daily, Review y Retrospectiva.
    \item \textbf{Control de Flujo Kanban \cite{anderson2010}:} Aporta visibilidad en tiempo real mediante un tablero descompuesto en fases, límites de trabajo en progreso (WIP) y gestión explícita de bloqueos.
    \item \textbf{Prácticas de XP \cite{beck2000}:} Garantiza la calidad técnica mediante integración continua, pruebas automatizadas y revisión de código entre pares.
\end{itemize}

\section{Prácticas Adoptadas del Proceso}
\subsection{Aportes de Scrum}
Sprints de 2 semanas con capacidad fija de 4 Historias de Usuario por iteración (1 HU asignada de principio a fin a cada desarrollador). Ceremonias obligatorias: Sprint Planning, Daily Scrum de 15 minutos, Sprint Review y Retrospectiva.

\section{Prácticas Adoptadas del Proceso}
\subsection{Aportes de Scrum}
Sprints de 2 semanas con cadencia regular para completar las 15 historias del MVP. Ceremonias formales: Sprint Planning con estimación en Planning Poker, Daily Scrum de 15 minutos, Sprint Review y Retrospectiva al término de cada iteración.

\subsection{Aportes de Kanban y Límites WIP}
\begin{itemize}[leftmargin=*]
    \item \textbf{Tablero Visual de Cuatro Columnas:} Estructura controlada: \texttt{TO DO} $\rightarrow$ \texttt{DESARROLLO [WIP 3]} $\rightarrow$ \texttt{REVISIÓN [WIP 2]} $\rightarrow$ \texttt{DONE}.
    \item \textbf{Límites WIP Justificados:} Para los 3 desarrolladores activos, se establece $\text{WIP} = 3$ en Desarrollo (1 desarrollador = máximo 1 HU activa) y $\text{WIP} = 2$ en Revisión.
    \item \textbf{Principio de Fluidez:} Al alcanzarse la capacidad de Revisión (2/2), la prioridad del equipo cambia inmediatamente: suspender el inicio de nuevo trabajo y colaborar para finalizar y fusionar las historias en revisión (\textit{``Stop starting, start finishing''}).
    \item \textbf{Visibilidad de Impedimentos:} Bloqueos gestionados mediante etiqueta \texttt{BLOQUEADA} en la columna actual, sin crear columnas redundantes.
\end{itemize}

\subsection{Aportes de Extreme Programming (XP) y Ciclo TDD}
Cada Historia de Usuario desglosa internamente cuatro subtareas técnicas bajo el ciclo de Desarrollo Guiado por Pruebas:
\begin{enumerate}[leftmargin=*]
    \item \textbf{T1 -- Pruebas / RED:} Escritura inicial de pruebas unitarias o de integración que fallan.
    \item \textbf{T2 -- Implementación / GREEN:} Codificación mínima necesaria para superar las pruebas.
    \item \textbf{T3 -- Refactorización:} Optimización de estructura y legibilidad sin alterar el comportamiento.
    \item \textbf{T4 -- Validación e Integración:} Verificación de criterios BDD (\textit{Dado/Cuando/Entonces}) y preparación del Pull Request.
\end{enumerate}
Complementado con Integración Continua (CI), revisión obligatoria por pares (\textit{Peer Review}) y diseño simple.

\subsection{Uso de Inteligencia Artificial en el Proceso}
Uso de modelos de lenguaje para asistencia en generación de fixtures, casos de prueba y esquemas Prisma, bajo el principio de \textbf{validación humana obligatoria}: ningún miembro incorpora código sin comprenderlo y probarlo rigurosamente.

\section{Roles, Responsabilidades y Forma de Coordinación}
\subsection{Política de Roles Rotativos por Sprint}
El equipo adopta una estructura de \textbf{roles rotativos por iteración} para fomentar el aprendizaje integral de todos los integrantes:
\begin{itemize}[leftmargin=*]
    \item \textbf{Sprint 1 (Sprint Piloto de Evaluación):}
    \begin{itemize}
        \item \textbf{Product Owner:} Saire Bustamante, Edmil Jampier. Responsable de la priorización del backlog, validación de criterios BDD (\textit{Dado/Cuando/Entonces}) y desarrollo full-stack.
        \item \textbf{Scrum Master:} Barazorda Cuellar, Hector. Responsable de facilitar ceremonias (Daily de 15 min), remoción de bloqueos técnicos y desarrollo full-stack.
        \item \textbf{Desarrolladores Full-Stack:} Carpio Hermoza, Alex; Quispe Mamani, Domingo de Guzman; Ticona Jancco, Ronaldo. Ejecución autónoma de sus HUs asignadas y revisión por pares.
    \end{itemize}
    \item \textbf{Evaluación y Rotación en Retrospectiva:} Al cierre del Sprint 1, el equipo analiza la dinámica y efectividad de las acciones para determinar la rotación de los roles de Product Owner y Scrum Master hacia los demás integrantes en los Sprints 2 y 3.
\end{itemize}

\subsection{Estrategia de Ramas en Git y Flujo para Desarrolladores}
Para garantizar orden en el código y evitar colisiones, se define el siguiente esquema de ramas:
\begin{itemize}[leftmargin=*]
    \item \textbf{Rama Principal (\texttt{edmil-saire}):} Rama \textit{main} del repositorio oficial, destinada exclusivamente a versiones estables y libres de errores.
    \item \textbf{Rama de Integración (\texttt{desarrollo}):} Espacio común donde convergen los avances para testear la unión de los desarrolladores antes de producción.
    \item \textbf{Ramas Individuales por Desarrollador:} Cada miembro dispone de una rama con su nombre (\texttt{hector-barazorda}, \texttt{alex-carpio}, \texttt{domingo-quispe}, \texttt{edmil-saire}, \texttt{ronaldo-ticona}).
    \item \textbf{Flujo Operativo:} Cada desarrollador ejecuta \texttt{git pull origin desarrollo} en su rama para sincronizar cambios, sube commits con su cuenta de GitHub a su rama personal y genera un Pull Request hacia \texttt{desarrollo}. Tras el testeo conjunto, se integra a \texttt{edmil-saire}.
\end{itemize}

\section{Flujo de Trabajo del Equipo}
\subsection{Estructura del Tablero Kanban Oficial y Límites WIP}
El flujo se organiza formalmente en cuatro estados con límites explícitos de trabajo en progreso:

\begin{table}[H]
\centering
\small
\begin{tabularx}{\linewidth}{|l|X|c|}
\hline
\rowcolor{gray!15}
\textbf{Columna} & \textbf{Significado Operativo} & \textbf{Límite WIP} \\
\hline
\textbf{TO DO} & Historias del Sprint comprometidas aún no iniciadas. & Sin límite \\
\hline
\textbf{DESARROLLO} & Historias en codificación activa y construcción de pruebas. & \textbf{3} \\
\hline
\textbf{REVISIÓN} & Desarrollo culminado; en espera de revisión por pares y testing de integración. & \textbf{2} \\
\hline
\textbf{DONE} & Historias integradas en \texttt{desarrollo} y aceptadas por el Product Owner. & Sin límite \\
\hline
\end{tabularx}
\caption{Estructura de columnas y políticas WIP del tablero Kanban.}
\label{tab:kanban}
\end{table}

\subsection{Justificación de los Límites WIP para 3 Desarrolladores}
\begin{itemize}[leftmargin=*]
    \item \textbf{Desarrollo = 3:} Dado que el equipo cuenta con 3 integrantes en el rol Development, la regla es estricta: \textbf{1 desarrollador = máximo 1 HU activa en Desarrollo}. Esto evita la fragmentación de atención y asegura culminación vertical antes de iniciar otra tarea.
    \item \textbf{Revisión = 2:} Previene la formación de cuellos de botella en la integración. Si la columna se satura (2/2), ningún desarrollador puede enviar más código a Revisión; la prioridad absoluta del equipo pasa a ser colaborar en revisar y fusionar las historias pendientes.
\end{itemize}

\subsection{Unidad de Flujo y Subtareas Técnicas TDD}
La tarjeta que transita por el tablero es la \textbf{Historia de Usuario} completa. Internamente en Jira, cada HU desglosa las 4 subtareas técnicas del ciclo TDD:
\begin{center}
\texttt{[T1: Pruebas / RED]} $\rightarrow$ \texttt{[T2: Implementación / GREEN]} $\rightarrow$ \texttt{[T3: Refactorización]} $\rightarrow$ \texttt{[T4: Validación e Integración]}
\end{center}

\subsection{Políticas Explícitas de Transición}
\begin{itemize}[leftmargin=*]
    \item \textbf{TO DO $\rightarrow$ DESARROLLO:} HU perteneciente al sprint, con criterios BDD definidos, desarrollador asignado y espacio en WIP de Desarrollo ($\le 3$).
    \item \textbf{DESARROLLO $\rightarrow$ REVISIÓN:} Subtareas $T1$, $T2$ y $T3$ completadas, pruebas locales en verde, commits en la rama individual y espacio en WIP de Revisión ($\le 2$).
    \item \textbf{REVISIÓN $\rightarrow$ DONE:} Aprobación en \textit{Peer Review}, integración sin conflictos en \texttt{desarrollo}, criterios de aceptación verificados y validación por el Product Owner.
\end{itemize}

\subsection{Gestión de Bloqueos e Impedimentos}
No se crea una columna independiente para bloqueos. La tarjeta permanece en su estado actual con la etiqueta \texttt{BLOQUEADA}. Se registra en Jira el motivo técnico, el responsable y la acción de mitigación, interviniendo el Scrum Master en la Daily Scrum para su resolución prioritaria.

\subsection{Dinámica por Sprints}
\begin{itemize}[leftmargin=*]
    \item \textbf{Sprint 1 (HU-001 a HU-006):} Arranque con HU-001 (Dev 1), HU-002 (Dev 2) y HU-003 (Dev 3) en Desarrollo. Al culminar, ingresan ordenadamente HU-004, HU-005 y HU-006 respetando el límite de 3.
    \item \textbf{Sprint 2 (HU-007 a HU-012):} Arranque con HU-007, HU-008 y HU-009 en Desarrollo; HU-010, HU-011 y HU-012 en To Do con ingreso progresivo.
    \item \textbf{Sprint 3 (HU-013 a HU-015):} Coincidencia de 3 HUs para 3 Desarrolladores. Gestión de dependencias: HU-014 coordina su integración tras la definición de datos de HU-013.
\end{itemize}

\subsection{Definición de Hecho (Definition of Done)}
Una HU es declarada \textit{Done} si supera el 100\% de criterios de aceptación BDD, cuenta con pruebas automáticas en verde, no presenta errores de TypeScript/Next.js, tiene persistencia validada en MongoDB Atlas con Prisma, y se encuentra fusionada en la rama \texttt{desarrollo}.

\section{Incrementos Presentados del Proyecto}
\subsection{Incremento 1 (Sprint 1: Núcleo de Capacitación y Acceso Público)}
\begin{itemize}[leftmargin=*]
    \item \textbf{HU-001 (Acceder al área interna):} Autenticación segura y redirección por rol (Dev 1: Carpio Hermoza, Alex).
    \item \textbf{HU-002 (Crear capacitación):} Gestión de datos generales y persistencia en MongoDB (Dev 2: Quispe Mamani, Domingo).
    \item \textbf{HU-003 (Slides interactivos):} Constructor de diapositivas con texto, imágenes y botones (Dev 3: Ticona Jancco, Ronaldo).
    \item \textbf{HU-004 (Publicar capacitación):} Previsualización interactiva y flujo de estados (Dev 1: Carpio Hermoza, Alex).
    \item \textbf{HU-005 (Explorar capacitaciones públicas):} Catálogo abierto para estudiantes y docentes (Dev 2: Quispe Mamani, Domingo).
    \item \textbf{HU-006 (Consultar capacitación pública):} Visualizador de slides interactivo sin autenticación (Dev 3: Ticona Jancco, Ronaldo).
\end{itemize}
\textbf{Pruebas y Evidencias:} Pruebas automatizadas sobre endpoints de Next.js, verificación de persistencia en MongoDB Atlas Cluster0 vía Prisma ORM y validación de los criterios CA-01 a CA-04 de cada historia bajo TDD.

\section{Métricas y Seguimiento del Proceso}
\begin{itemize}[leftmargin=*]
    \item \textbf{Velocidad del Sprint:} 6 HUs en Sprint 1, 6 HUs en Sprint 2 y 3 HUs en Sprint 3 (Total: 15 Historias de Usuario del MVP completadas).
    \item \textbf{Cycle Time:} Promedio de 3 a 4 días hábiles por historia desde \texttt{TO DO} hasta \texttt{DONE}.
    \item \textbf{Criterios Superados:} Tasa de éxito del 100\% en criterios de aceptación verificados.
    \item \textbf{Calidad de Código:} Cero regresiones en la rama \texttt{desarrollo} respaldadas por la suite TDD.
\end{itemize}

\section{Acuerdos de Mejora para el Siguiente Incremento}
\begin{enumerate}[leftmargin=*]
    \item Aprobar prototipos en Figma a más tardar el segundo día de cada sprint.
    \item Fortalecer los tests de integración con Supertest para casos de sesión inválida.
    \item Coordinar el esquema de base de datos relacional del Sprint 2 durante la sesión de Sprint Planning.
\end{enumerate}

\section{Política de Uso de IA y Registro de Uso}
\begin{table}[H]
\centering
\small
\begin{tabularx}{\linewidth}{|l|l|X|X|}
\hline
\rowcolor{gray!15}
\textbf{Fecha} & \textbf{Herramienta} & \textbf{Actividad / Propósito} & \textbf{Validación Técnica} \\
\hline
30/09/2026 & Gemini / Claude & Estructuración formal de la especificación técnica CGB Academy & Revisión contra el documento base oficial de 18 páginas. \\
\hline
30/09/2026 & Copilot / ChatGPT & Formateo y tipografía del informe en LaTeX (\texttt{informe.tex}) & Compilación limpia con \texttt{pdflatex} y verificación de tablas. \\
\hline
30/09/2026 & Antigravity Assistant & Automatización de scripts de verificación y ramas Git & Ejecución local y verificación de sincronización en GitHub. \\
\hline
\end{tabularx}
\caption{Bitácora de uso asistido de Inteligencia Artificial.}
\label{tab:ia}
\end{table}

\section{Conclusiones}
\begin{enumerate}[leftmargin=*]
    \item El método híbrido garantiza que cada entrega sea un incremento vertical completo con diseño, backend, frontend y pruebas.
    \item La delimitación estricta del MVP de CGB Academy previene la dispersión del esfuerzo técnico y asegura la entrega a tiempo.
    \item La adopción ética de IA incrementa sustancialmente la productividad sin comprometer el rigor ingenieril ni el aprendizaje formativo.
\end{enumerate}

\bibliographystyle{plain}
\bibliography{referencias}

\end{document}
